\documentclass[11pt]{article}
\usepackage[T1]{fontenc}
\usepackage[utf8]{inputenc}
\usepackage[french]{babel}
\usepackage{amsmath,amssymb,amsthm}
\usepackage{hyperref}
\setlength{\emergencystretch}{1em}
\newcommand{\End}{\operatorname{End}}
\newcommand{\ev}{\operatorname{ev}}
\title{Cayley--Hamilton et transport par univalence}
\author{}
\date{27 octobre 2026}
\begin{document}
\maketitle

\section{Données et conventions}
Soit $A$ un anneau commutatif dont le porteur est un ensemble, et $M$
un $A$-module. L'algèbre $\End_A(M)$ a pour éléments les applications
linéaires, avec addition et action scalaire ponctuelles, produit
$fg=f\circ g$ et unité $\mathrm{id}_M$. Cette construction est indépendante
de toute base. Les lois de linéarité sont des propositions : l'égalité
ponctuelle des applications suffit pour l'égalité des éléments de l'algèbre.

Les matrices $B$ sont stockées par lignes. Leur action est
$(Bv)_i=\sum_j B_{ij}v_j$. La bibliothèque de déterminants stocke les
colonnes : le même tableau y est fourni par $a(j)(i)=B_{ij}$.
Ainsi $\chi(B)=\det(TI-B)$ et le théorème matriciel existant donne
$\ev_{M_n(A)}(B,\chi(B))=0$. Aucune hypothèse $1\ne0$ n'est nécessaire.

\section{Assemblage des algèbres}
L'application $B\mapsto(v\mapsto Bv)$ est un isomorphisme
$M_n(A)\simeq\End_A(A^n)$. Son inverse prend les images des vecteurs
de base pour colonnes ; l'identité $v=\sum_jv_je_j$ prouve que ces
deux applications sont inverses. Le produit correspond à la composition
par échange de deux sommes finies.

Une équivalence linéaire $e:A^n\simeq M$ donne ensuite
\[
 \Phi_e(B)=e\circ\operatorname{action}(B)\circ e^{-1}.
\]
Elle préserve les opérations d'algèbre ; la compatibilité au produit
utilise $e^{-1}e=\mathrm{id}$. Le principe d'identité des structures donne
un chemin $s_e:M_n(A)=\End_A(M)$.

\section{Le calcul du transport}
Pour des algèbres $E,F$ et $\phi:E\simeq F$, le chemin $s=\mathrm{sip}(\phi)$
satisfait $\mathrm{idToIso}(s)=\phi$. Par induction sur un chemin,
le transport dans la famille $E\mapsto E.A$ est l'application sous-jacente
à $\mathrm{idToIso}(s)$. Il vaut donc $\phi$ pour $s=\mathrm{sip}(\phi)$.
Les mêmes inductions donnent
\[
 \operatorname{transport}_s(\ev_E(t,p))=\ev_F(\phi(t),p).
\]
Le polynôme $p\in A[T]$ reste fixe. Pour transporter une identité
d'annulation, on utilise le type des algèbres pointées $(E,t)$ et la
famille $\ev_E(t,p)=0_E$ sur ce type. Il ne s'agit pas d'une réduction
définitionnelle de l'axiome d'univalence.

Plus généralement on transporte le paquet
\[
 \sum_{c:E.A\to A[T]}\ \prod_{t:E.A}\ev_E(t,c(t))=0_E.
\]
La fonction du paquet transporté vaut $c\circ\phi^{-1}$ ; cette formule
se démontre par induction sur les chemins, puis par le calcul de SIP.
Pour $E=M_n(A)$ on obtient, pour tout $f\in\End_A(M)$,
\[
 \ev_{\End_A(M)}(f,\chi(\Phi_e^{-1}(f)))=0.
\]
On peut évaluer cette égalité sur tout vecteur de $M$.

\section{Comparaison intrinsèque : second jalon}
Le paquet précédent dépend a priori de $e$. L'univalence ne démontre
pas à elle seule que la fonction caractéristique transportée ne dépend
pas de cette base. Un polynôme annulateur monique n'est pas unique.

Une définition intrinsèque usuelle est
\[
 \chi_M(f)=\det_{A[T]}(T\mathrm{id}_{M[T]}-f[T]),
 \qquad M[T]=A[T]\otimes_A M.
\]
Le déterminant est le scalaire par lequel l'endomorphisme induit agit
sur la puissance extérieure supérieure. Le module libre de rang $n$
peut n'avoir de base que sous troncation : existence et unicité du
scalaire permettent de définir ce déterminant sans choisir de base.

Pour éviter de construire une nouvelle bibliothèque de changement de base,
la formalisation utilise une variante intrinsèque équivalente en coordonnées :
\[
 \begin{gathered}
 D(M)=\operatorname{Hom}_A(M,A[T]),\qquad f^*(\lambda)=\lambda\circ f,\\
 \chi_M(f)=\det_{A[T]}(T\mathrm{id}_{D(M)}-f^*).
 \end{gathered}
\]
Le module $D(M)$ est défini sans base : les opérations et la multiplication
par un polynôme sont ponctuelles. La précomposition par $f$ est
$A[T]$-linéaire. Elle renverse l'ordre des produits ; on ne l'utilise
donc pas comme un isomorphisme d'algèbres d'endomorphismes.

Une base, écrite cette fois $e:M\simeq A^n$, fournit les vecteurs
$v_j=e^{-1}(e_j)$ et l'équivalence
\[
 D(M)\longrightarrow A[T]^n,\quad
 \lambda\longmapsto(\lambda(v_j))_j.
\]
Son inverse envoie $(p_j)_j$ sur
$x\mapsto\sum_j\iota(e(x)_j)p_j$, où $\iota:A\to A[T]$ est l'inclusion
des constantes. L'identité $x=\sum_je(x)_jv_j$ démontre que ces fonctions
sont inverses. Si la base n'existe que sous troncation, cette construction
donne une base de $D(M)$ sous troncation, ce qui suffit à définir son
déterminant intrinsèque avec \texttt{DetEnd}.

Si $a_{ij}$ est la matrice de $f$ (indices usuels de lignes et colonnes),
la matrice de $T\mathrm{id}-f^*$ est $(TI-a)^t$. La formule
\texttt{DetEnd.delta-coords}, puis l'invariance du déterminant par
transposition, donnent
\[
 \chi_M(f)=\det((TI-a)^t)=\det(TI-a).
\]
On obtient ainsi indépendance de la base, monicité au sens des coefficients
et Cayley--Hamilton intrinsèque par comparaison avec le paquet transporté
du premier jalon. L'élimination de la base sous troncation se fait dans
l'identité d'annulation, qui est une proposition.

Pour une équivalence linéaire $j:M\simeq N$, transporter une base de $M$
en une base de $N$ donne la même matrice pour $f$ et $jfj^{-1}$.
La comparaison précédente prouve
$\chi_N(jfj^{-1})=\chi_M(f)$, également sans choix de base.
La comparaison avec une construction tensorielle séparée de $M[T]$ n'est
pas formalisée ; elle n'est pas nécessaire aux résultats ci-dessus.

\section{Résultats et exemples vérifiés}
Les deux jalons sont formalisés dans \texttt{sag/ch-univalence/}.
Le client \texttt{EndCayleyHamilton} transporte effectivement un paquet
par \texttt{AlgSIP.sip}. Le théorème intrinsèque final est :
\begin{quote}\small\ttfamily
CharPolyComparison.Over.At.cayley-hamilton
\end{quote}
Ses dépendances passent
par le théorème matriciel existant, le transport du paquet, puis la
comparaison du déterminant intrinsèque. Aucun nouvel axiome n'est ajouté
au contexte \texttt{Theory} des bibliothèques.

Les exemples vérifient les rangs zéro et un, puis la matrice symbolique
non symétrique $\left(\begin{smallmatrix}u&v\\w&z\end{smallmatrix}\right)$.
Le polynôme est identifié à
\[
 (uz-vw)+T(-(u+z)+T),
\]
et, sur tout vecteur $x$, le théorème transporté donne
\[
 f(f(x))-(u+z)f(x)+(uz-vw)x=0.
\]
Un cisaillement $s_c(x,y)=(x+cy,y)$, d'inverse $s_{-c}$, vérifie aussi
la formule explicite de conjugaison et l'invariance du polynôme intrinsèque.

\section{Ce que mesure l'expérience}
Le transport d'un élément, d'une évaluation et d'un paquet est général :
il n'utilise ni les coefficients de la matrice ni une induction sur les
polynômes. Les seules identifications propres au client sont l'isomorphisme
matrices--endomorphismes et la formule de coordonnées. Dans le second
jalon, l'identification de la matrice obtenue par les équivalences inverses
avec celle du déterminant se vérifie par réflexivité : c'est une conversion
entre les coordonnées explicites, et non une réduction de l'univalence.

Le calcul de \texttt{sip} est utilisé exclusivement par ses lemmes publics,
sans dérouler sa preuve opaque. Le petit client et les lemmes généraux
montrent ici une réutilisation effective. Les temps et le nombre exact
de déclarations figurent dans le \texttt{README.txt}. Ils ne permettent
pas à eux seuls de conclure à la tenue du système pour une bibliothèque
de la taille d'une formalisation de FLT.
\end{document}
